% Tapis de Quillen, with Quillen's offprint and preprint (cote 162-5, typed part): archive cote 162-5, site page 7 (PDF page 8).
% typescript. Draft compared: none.
% AI-transcribed by Claude Opus 5.5 (claude-opus-5-5, adaptive thinking, effort high), Message Batches API, batch t162_full (msgbatch_01UAh3zyWPd454AeDkhnmQ33), 2026-11-13; not yet checked by eye.
\scanpage{162-5/7}
\mspage{\hand{6}}

morphismes de flêches.

Notons une symmétrie évidente dans la c atégorie (Bicat), interprété
en termes d\xout{e}'ensemble bi-semisimpliciaux, obtenue en échangeant les lign
et les colonnes dqns \xout{\illegible{1}} de tels ensembles bi... Un ensemble bisemisimpl
al qui ... peut donc de deux façons différentes être considéré comme déf
nissant une bicatégorie. De même le groupe symétrique $\mathrm{Sym}_n$ opère sur
la catégorie \xout{de} (nupcat) des catégories n-uples ..

Une 2-catégorie peut être définie comme constituée par un ensemble
\xout{discret} $C_0^{\hand{0}}$ (l'ensemble des objets ), qu'on interprètera comme une catégo
rie discrète, plus la donnée pour tout $x,y \hand{$\in$} C_0^{\hand{0}}$ d'une catégorie $\underline{\mathrm{Hom}}(x,y)$,
plus \xout{\illegible{1}} \inserted[typed]{des foncteurs} composition entre celles-ci ... Prenant la catégorie somme \xout{\illegible{1}} C
des $\underline{\mathrm{Hom}}(x,y)$, On peut exprimer la donnée de ces foncteurs de composition
et des sources et buts\inserted{\unsure{$x,y$ )}} dans les catégories $\underline{\mathrm{Hom}}(x,y)$, comme définis par
des foncteurs "source" et "but" $C_{1}^{\hand{1}} \to C_{0}^{\hand{0}}$, et un foncteur de composition
$C_{1}^{\hand{1}} \times_{C_0^{\hand{0}}} C_{1}^{\hand{1}} \longrightarrow C_1^{\hand{1}}$ , les axiomes des 2-catégories correspondant exactement au
axiomes des(Cat)-catégories. De cette façon, la catégorie des 2-catégor
est équivalente à la catégorie des (Cat)-catégories qui ont une "catégo
d'objets" qui est discrète, ou encore à la catégorie des Bicatégories
qui, interprétées comme ensembles bi-semisimpliciaux, ont \xout{une} comme
première ligne un ensemble semisimplicial discret (i.e. un foncteur cons
tant sur \hand{$\Delta$} ). Il parait qu'on a une interprétation simple dans le même
esprit de la notion de n-catégorie en termes d'ensembles n-semi-simplici
\margin{$\|$} \xout{au} Demander à Quillen relations entre la notion de m-Catégorie dans
\margin{$\|$} la catégorie (n-Cat), et celle d\over{e}{'} ensemble (m+n)-semisimplicial.
\margin{$\|$} \xout{(en particulier entre cette dernière notion et celle de}

L'interprétation standard semisimpliciale permet d'interpréter

\margin{7}
